\begin{proposition}[pullback tower; Kahn I, \S\S7.3.1--7.3.2, pp.~38--43]\label{prop:kahn-pullback-tower}
\begin{enumerate}
\item (Existence; Lemma~7.8, p.~40.)  In the category of complex
\(1\)-manifolds, the pullback of a diagram \(g\colon G\to V\),
\(h\colon H\to V\) exists whenever \(g\) (or, symmetrically, \(h\)) is
an immersion; its underlying topological space is the topological
fiber product
\(E=\{(x,y)\in G\times H:\ g(x)=h(y)\}\)
with the coordinate projections \(\pi_G,\pi_H\) (Kahn describes \(E\)
as the \emph{equalizer}; \S7.3.1, p.~39).
\item (Inheritance; Lemma~7.7, p.~40 and Lemma~7.9, p.~40.)  Each projection inherits
the properties of the opposite leg: if \(g\) is proper (resp.\ a local
homeomorphism, a homeomorphism, injective, surjective, bijective),
then so is \(\pi_H\); and for \(y\in H\) the fiber \(\pi_H^{-1}(y)\)
is in natural bijection with \(g^{-1}(h(y))\).  If \(g\) is an
immersion and \(h\) is proper of degree \(d\), then \(\pi_H\) is an
immersion and \(\pi_G\) is proper of degree \(d\), and the local
degree of \(\pi_G\) at \(z\in E\) equals the local degree of \(h\) at
\(\pi_H(z)\).
\item (One \(\psi\)-polynomial-like pullback; \S7.3.2, p.~42 and Lemma~7.11, p.~42.)  For a
degree-\(d\) \(\psi\)-polynomial-like map
\((i,f)\colon(\BU',K')\to(\BU,K)\), the pullback of
\(i\colon\BU'\to\BU\) along \(f\colon\BU'\to\BU\) is
\[
 \BU''=\{(x,y)\in\BU'\times\BU':\ i(x)=f(y)\},\qquad
 i_1(x,y)=y,\quad f_1(x,y)=x,
\]
and \((i_1,f_1)\colon(\BU'',K'')\to(\BU',K')\) is again a
degree-\(d\) \(\psi\)-polynomial-like map, where
\(K''=i_1^{-1}(K')=f_1^{-1}(K')\).
\item (Tower; \S7.3.2, p.~43.)  Relabeling \((i,f)\) as
\((i_0,f_0)\colon(\BU^1,K^1)\to(\BU^0,K^0)\) and iterating produces
degree-\(d\) \(\psi\)-polynomial-like maps
\((i_n,f_n)\colon(\BU^{n+1},K^{n+1})\to(\BU^n,K^n)\) with the
commuting squares
\[
 i_n\circ f_{n+1}=f_n\circ i_{n+1}\qquad(n\in\N),
\]
each square being a genuine categorical pullback (fiber product) of
the previous stage.
\item (Germ preservation; \S7.3.3, p.~43.)  If \(U'\subset\BU'\),
\(U\subset\BU\) are such that \((i,f)\colon U'\to U\) is
polynomial-like (\(i(U')\Subset U\) and \(i|_{U'}\) an embedding),
then the pullback \(U''\) of \((i,f)\colon U'\to U\) sits inside
\(\BU''\) and \((i_1,f_1)\colon U''\to U'\) is polynomial-like;
consequently the polynomial-like germ of \((i_1,f_1)\) equals that of
\((i,f)\).
\end{enumerate}
\end{proposition}

\begin{proof}[Proof (faithful paraphrase of Kahn I, \S7.3.1, pp.~40--41 and \S7.3.2, p.~42)]
(1) By the inheritance statement for local homeomorphisms,
\(\pi_H\colon E\to H\) is a local homeomorphism when \(g\) is an
immersion; pulling the charts of \(H\) back through \(\pi_H\) endows
\(E\) with a complex structure for which \(\pi_H\) is holomorphic, and
locally \(\pi_G=g^{-1}\circ h\circ\pi_H\), hence \(\pi_G\) is
holomorphic.  The universal property is checked by factoring any
competitor \(F\to E\) through the topological pullback and verifying
holomorphy locally via \((\pi_H|_{\widehat U})^{-1}\circ f_H\)
(proof of Lemma~7.8, pp.~40--41).  Kahn notes the pullback may be disconnected, which is
why he writes ``complex \(1\)-manifold'' rather than ``Riemann
surface'' (\S7.3.1, p.~41).

(2) Properness: for compact \(B\subset H\), the set \(\pi_H^{-1}(B)\)
is a closed subset of \(g^{-1}(h(B))\times B\), hence compact
(proof of Lemma~7.7, p.~40).  Local homeomorphism: over a neighborhood on which
\(g\) is invertible, \(\pi_H\) has the explicit inverse
\(y\mapsto((g|_U)^{-1}(h(y)),y)\) (ibid.).  The set-theoretic
statements and the fiber correspondence are elementary
(ibid.).  When \(h\) is proper of degree \(d\), evaluating the
fiber correspondence at a regular value gives
\(|\pi_G^{-1}(x)|=d\), and the local-degree statement follows by
comparing small neighborhoods (proof of Lemma~7.9, pp.~40--41).

(3) By (1)--(2), \(i_1\) is an immersion and \(f_1\) is proper of
degree \(d\).  For \((x,y)\in\BU''\), the chain
\(x\in K'\iff i(x)\in K\iff f(y)\in K\iff y\in K'\) shows
\(i_1^{-1}(K')=f_1^{-1}(K')=:K''\) (\S7.3.2, p.~42).  Since
\(i\colon K'\to K\) is a homeomorphism, so is
\(i_1\colon K''\to K'\) (inheritance), whence \(K''\) is connected
(\S7.3.2, p.~42).  Every component of \(\BU''\) meets \(K''\) because
\(f_1\) is proper onto \(\BU'\); hence \(\BU''\) is connected
(\S7.3.2, p.~42).  A point of \(\BU''\) is critical for \(f_1\) exactly
when its \(i_1\)-image is critical for \(f\), with the same
multiplicity; therefore all critical points of \(f_1\) lie in \(K''\)
and their total multiplicity is \(d-1\), so Riemann--Hurwitz gives
\(\chi(\BU'')=1\) and \(\BU''\) is a disk (\S7.3.2, p.~42).  This
verifies all clauses of the \(\psi\)-polynomial-like definition
(Kahn~I, Lemma~7.11, p.~42).

(4) Apply (3) inductively; the commuting relation is the defining
property of the pullback square at each stage (\S7.3.2, p.~43).

(5) The pullback of the polynomial-like restriction embeds in the
pullback of the ambient pair; the restricted pair
\((i_1,f_1)\colon U''\to U'\) satisfies the polynomial-like
conditions, and its germ at the filled Julia set coincides with that
of \((i,f)\) (\S7.3.3, p.~43).
\end{proof}